\documentclass[10pt,a4paper]{amsart}
\usepackage{amsmath,amssymb,amsthm,graphicx}
\usepackage{verbatim}
\newenvironment{reference}{\comment}{\endcomment}
%\newenvironment{reference}{}{}
\newenvironment{slogan}{\comment}{\endcomment}
\newenvironment{history}{\comment}{\endcomment}
% For commutative diagrams we use Xy-pic
\usepackage[all]{xy}

% We use 2cell for 2-commutative diagrams.
\xyoption{2cell}
% Theorem environments.
%
\theoremstyle{plain}
\newtheorem{theorem}[subsection]{Theorem}
\newtheorem{proposition}[subsection]{Proposition}
\newtheorem{lemma}[subsection]{Lemma}

\theoremstyle{definition}
\newtheorem{definition}[subsection]{Definition}
\newtheorem{example}[subsection]{Example}
\newtheorem{exercise}[subsection]{Exercise}
\newtheorem{situation}[subsection]{Situation}

\theoremstyle{remark}
\newtheorem{remark}[subsection]{Remark}
\newtheorem{remarks}[subsection]{Remarks}

\numberwithin{equation}{subsection}

% Macros
%
\def\lim{\mathop{\mathrm{lim}}\nolimits}
\def\colim{\mathop{\mathrm{colim}}\nolimits}
\def\Spec{\mathop{\mathrm{Spec}}}
\def\Hom{\mathop{\mathrm{Hom}}\nolimits}
\def\Ext{\mathop{\mathrm{Ext}}\nolimits}
\def\SheafHom{\mathop{\mathcal{H}\!\mathit{om}}\nolimits}
\def\SheafExt{\mathop{\mathcal{E}\!\mathit{xt}}\nolimits}
\def\Sch{\mathit{Sch}}
\def\Mor{\mathop{\mathrm{Mor}}\nolimits}
\def\Ob{\mathop{\mathrm{Ob}}\nolimits}
\def\Sh{\mathop{\mathit{Sh}}\nolimits}
\def\NL{\mathop{N\!L}\nolimits}
\def\CH{\mathop{\mathrm{CH}}\nolimits}
\def\proetale{{pro\text{-}\acute{e}tale}}
\def\etale{{\acute{e}tale}}
\def\QCoh{\mathit{QCoh}}
\def\Ker{\mathop{\mathrm{Ker}}}
\def\Im{\mathop{\mathrm{Im}}}
\def\Coker{\mathop{\mathrm{Coker}}}
\def\Coim{\mathop{\mathrm{Coim}}}

% Boxtimes
%
\DeclareMathSymbol{\boxtimes}{\mathbin}{AMSa}{"02}

%
% Macros for moduli stacks/spaces
%
\def\QCohstack{\mathcal{QC}\!\mathit{oh}}
\def\Cohstack{\mathcal{C}\!\mathit{oh}}
\def\Spacesstack{\mathcal{S}\!\mathit{paces}}
\def\Quotfunctor{\mathrm{Quot}}
\def\Hilbfunctor{\mathrm{Hilb}}
\def\Curvesstack{\mathcal{C}\!\mathit{urves}}
\def\Polarizedstack{\mathcal{P}\!\mathit{olarized}}
\def\Complexesstack{\mathcal{C}\!\mathit{omplexes}}
% \Pic is the operator that assigns to X its picard group, usage \Pic(X)
% \Picardstack_{X/B} denotes the Picard stack of X over B
% \Picardfunctor_{X/B} denotes the Picard functor of X over B
\def\Pic{\mathop{\mathrm{Pic}}\nolimits}
\def\Picardstack{\mathcal{P}\!\mathit{ic}}
\def\Picardfunctor{\mathrm{Pic}}
\def\Deformationcategory{\mathcal{D}\!\mathit{ef}}
\usepackage{cleveref}
\crefname{originallemma}{original lemma}{original lemmas}
\Crefname{originallemma}{Original lemma}{Original lemmas}
\makeatletter
\newlabel{lemma-characterize-quasi-DM}{{lemma-characterize-quasi-DM}{source}{}{sourcepointer-1}{}}
\newlabel{lemma-characterize-quasi-DM@cref}{{[originallemma][0][]lemma-characterize-quasi-DM}{[1][0][]source}{}{}{}}
\makeatother
\crefname{originallemma}{original lemma}{original lemmas}
\Crefname{originallemma}{Original lemma}{Original lemmas}
\makeatletter
\newlabel{lemma-check-finite-type-covering}{{lemma-check-finite-type-covering}{source}{}{sourcepointer-2}{}}
\newlabel{lemma-check-finite-type-covering@cref}{{[originallemma][0][]lemma-check-finite-type-covering}{[1][0][]source}{}{}{}}
\makeatother
\crefname{originallemma}{original lemma}{original lemmas}
\Crefname{originallemma}{Original lemma}{Original lemmas}
\makeatletter
\newlabel{lemma-check-finite-type-precompose}{{lemma-check-finite-type-precompose}{source}{}{sourcepointer-3}{}}
\newlabel{lemma-check-finite-type-precompose@cref}{{[originallemma][0][]lemma-check-finite-type-precompose}{[1][0][]source}{}{}{}}
\makeatother
\crefname{originallemma}{original lemma}{original lemmas}
\Crefname{originallemma}{Original lemma}{Original lemmas}
\makeatletter
\newlabel{lemma-composition-locally-quasi-finite}{{lemma-composition-locally-quasi-finite}{source}{}{sourcepointer-4}{}}
\newlabel{lemma-composition-locally-quasi-finite@cref}{{[originallemma][0][]lemma-composition-locally-quasi-finite}{[1][0][]source}{}{}{}}
\makeatother
\crefname{originalreference}{original reference}{original references}
\Crefname{originalreference}{Original reference}{Original references}
\makeatletter
\newlabel{spaces-morphisms-lemma-locally-quasi-finite}{{spaces-morphisms-lemma-locally-quasi-finite}{source}{}{sourcepointer-5}{}}
\newlabel{spaces-morphisms-lemma-locally-quasi-finite@cref}{{[originalreference][0][]spaces-morphisms-lemma-locally-quasi-finite}{[1][0][]source}{}{}{}}
\makeatother
\title{Four equivalent characterizations of local quasi-finiteness for a morphism of algebraic stacks}
\author{The Stacks Project Authors}
\date{stacks-project:stacks-morphisms.tex:local-snapshot-20261012; GNU FDL 1.2 or later}
\begin{document}
\maketitle
\noindent\textit{Editorial note.} All delivered mathematical text is the author's own, copied from the pinned Stacks Project chapter source sources/stacks-algebraic-stacks-morphisms/pinned-original-source.tex; no statement or proof step was rewritten, reordered or removed.\par\medskip
\noindent\textit{Editorial note.} Preamble provenance: the Stacks Project keeps one single shared preamble for every chapter; this card takes it directly from the staged sources/stacks-algebraic-stacks-morphisms/preamble.tex as three verbatim spans (lines 8-12, 14-18 and 154-217), so no derived file was written and nothing was added or redefined. The wrapper class and the base packages are supplied by the assembler.\par\medskip
\noindent\textit{Editorial note.} Bibliography provenance: neither the counted statement nor its proof contains a citation (this lane rescans the whole span before writing the card), so the card carries no bibliography and no entry was reconstructed.\par\medskip
\noindent\textit{Editorial note.} Selection provenance: this card was opened under the file-based control contract (CONTROL.md), with the mechanical pre-card gate run on its statement and proof interval before the card was written.\par\medskip
\noindent\textit{Allowed standard prerequisites (editorial).}\begin{itemize}
\item The characterization of quasi-DM morphisms through locally quasi-finite covers. Invoked inside the counted proof (source lines 5160-5170). Links local quasi-finiteness with quasi-DM-ness in condition (2).
\item The criterion for checking a property of a morphism after a flat locally finitely presented cover. Invoked inside the counted proof (source lines 5171-5190). Allows the property to be tested on the covering algebraic space of condition (3).
\item The permanence of local quasi-finiteness under composition. Invoked inside the counted proof (source lines 5191-5205). Combines the two morphisms in conditions (3) and (4).
\end{itemize}\medskip
\ifdefined\hypertarget\hypertarget{sourcepointer-1}{}\fi
\noindent\textit{Source reference.} Original lemma lemma-characterize-quasi-DM: sources/stacks-algebraic-stacks-morphisms/pinned-original-source.tex, source lines 5086--5086. Navigation pointer only: the counted text refers to the Stacks item lemma-characterize-quasi-DM, which is an existing labeled result of this same chapter and is not part of this card. The pointer carries the target's own label as its number and its definition line as its span; its content is not reproduced here.\par\medskip
\ifdefined\hypertarget\hypertarget{sourcepointer-2}{}\fi
\noindent\textit{Source reference.} Original lemma lemma-check-finite-type-covering: sources/stacks-algebraic-stacks-morphisms/pinned-original-source.tex, source lines 3554--3554. Navigation pointer only: the counted text refers to the Stacks item lemma-check-finite-type-covering, which is an existing labeled result of this same chapter and is not part of this card. The pointer carries the target's own label as its number and its definition line as its span; its content is not reproduced here.\par\medskip
\ifdefined\hypertarget\hypertarget{sourcepointer-3}{}\fi
\noindent\textit{Source reference.} Original lemma lemma-check-finite-type-precompose: sources/stacks-algebraic-stacks-morphisms/pinned-original-source.tex, source lines 3586--3586. Navigation pointer only: the counted text refers to the Stacks item lemma-check-finite-type-precompose, which is an existing labeled result of this same chapter and is not part of this card. The pointer carries the target's own label as its number and its definition line as its span; its content is not reproduced here.\par\medskip
\ifdefined\hypertarget\hypertarget{sourcepointer-4}{}\fi
\noindent\textit{Source reference.} Original lemma lemma-composition-locally-quasi-finite: sources/stacks-algebraic-stacks-morphisms/pinned-original-source.tex, source lines 5037--5037. Navigation pointer only: the counted text refers to the Stacks item lemma-composition-locally-quasi-finite, which is an existing labeled result of this same chapter and is not part of this card. The pointer carries the target's own label as its number and its definition line as its span; its content is not reproduced here.\par\medskip
\ifdefined\hypertarget\hypertarget{sourcepointer-5}{}\fi
\noindent\textit{Source reference.} Original reference spaces-morphisms-lemma-locally-quasi-finite: sources/stacks-algebraic-stacks-morphisms/pinned-original-source.tex, source lines 4931--4931. Navigation pointer only: the counted text refers to the Stacks item spaces-morphisms-lemma-locally-quasi-finite, which belongs to another chapter of the project and is not part of this card. The pointer carries the target's own label as its number because this chapter's staged snapshot does not contain the other chapter's tag table; its content is not reproduced here.\par\medskip
\begin{lemma}
\label{lemma-characterize-locally-quasi-finite}
Let $f : \mathcal{X} \to \mathcal{Y}$ be a morphism of algebraic stacks.
The following are equivalent
\begin{enumerate}
\item $f$ is locally quasi-finite,
\item $f$ is quasi-DM and for any morphism $V \to \mathcal{Y}$ with $V$
an algebraic space and any locally quasi-finite morphism
$U \to \mathcal{X} \times_\mathcal{Y} V$ where $U$ is an algebraic space
the morphism $U \to V$ is locally quasi-finite,
\item for any morphism $V \to \mathcal{Y}$ from an algebraic space $V$
there exists a surjective, flat, locally finitely presented, and locally
quasi-finite morphism $U \to \mathcal{X} \times_\mathcal{Y} V$ where
$U$ is an algebraic space such that $U \to V$ is locally quasi-finite,
\item there exists algebraic spaces $U$, $V$, a surjective, flat,
and locally of finite presentation morphism $V \to \mathcal{Y}$,
and a morphism $U \to \mathcal{X} \times_\mathcal{Y} V$ which
is surjective, flat, locally of finite presentation, and
locally quasi-finite such that $U \to V$ is locally quasi-finite.
\end{enumerate}
\end{lemma}

\begin{proof}
Assume (1). Then $f$ is quasi-DM by assumption. Let
$V \to \mathcal{Y}$ and $U \to \mathcal{X} \times_\mathcal{Y} V$
be as in (2). By
Lemma \ref{lemma-composition-locally-quasi-finite}
the composition $U \to \mathcal{X} \times_\mathcal{Y} V \to V$ is
locally quasi-finite. Thus (1) implies (2).

\medskip\noindent
Assume (2). Let $V \to \mathcal{Y}$ be as in (3). By
Lemma \ref{lemma-characterize-quasi-DM}
we can find an algebraic space $U$ and a surjective, flat, locally
finitely presented, locally quasi-finite morphism
$U \to \mathcal{X} \times_\mathcal{Y} V$. By (2) the composition
$U \to V$ is locally quasi-finite. Thus (2) implies (3).

\medskip\noindent
It is immediate that (3) implies (4).

\medskip\noindent
Assume (4). We will prove (1) holds, which finishes the proof. By
Lemma \ref{lemma-characterize-quasi-DM}
we see that $f$ is quasi-DM. To prove that $f$ is locally of finite type
it suffices to prove that $g : \mathcal{X} \times_\mathcal{Y} V \to V$ is
locally of finite type, see
Lemma \ref{lemma-check-finite-type-covering}.
Then it suffices to check that $g$ precomposed with
$h : U \to \mathcal{X} \times_\mathcal{Y} V$ is locally of finite type, see
Lemma \ref{lemma-check-finite-type-precompose}.
Since $g \circ h : U \to V$ was assumed to be locally quasi-finite
this holds, hence $f$ is locally of finite type.
Finally, let $k$ be a field and let $\Spec(k) \to \mathcal{Y}$
be a morphism. Then $V \times_\mathcal{Y} \Spec(k)$ is
a nonempty algebraic space which is locally of finite presentation
over $k$. Hence we can find a finite extension $k'/k$ and
a morphism $\Spec(k') \to V$ such that
$$
\xymatrix{
\Spec(k') \ar[r] \ar[d] & V \ar[d] \\
\Spec(k) \ar[r] & \mathcal{Y}
}
$$
commutes (details omitted). Then $\mathcal{X}_{k'} \to \mathcal{X}_k$
is representable (by schemes), surjective, and finite locally free. In
particular $|\mathcal{X}_{k'}| \to |\mathcal{X}_k|$ is surjective and open.
Thus it suffices to prove that $|\mathcal{X}_{k'}|$ is discrete. Since
$$
U \times_V \Spec(k') =
U \times_{\mathcal{X} \times_\mathcal{Y} V} \mathcal{X}_{k'}
$$
we see that $U \times_V \Spec(k') \to \mathcal{X}_{k'}$ is
surjective, flat, and locally of finite presentation (as a base change
of $U \to \mathcal{X} \times_\mathcal{Y} V$). Hence
$|U \times_V \Spec(k')| \to |\mathcal{X}_{k'}|$ is surjective and
open. Thus it suffices to show that $|U \times_V \Spec(k')|$ is
discrete. This follows from the fact that $U \to V$ is locally
quasi-finite (either by our definition above or from the original definition
for morphisms of algebraic spaces, via
Morphisms of Spaces, Lemma \ref{spaces-morphisms-lemma-locally-quasi-finite}).
\end{proof}

\end{document}
